\documentclass[10pt, aspectratio=169]{beamer}

% --- TEMA Y ESTILO ---
\usetheme{Madrid}
\usecolortheme{whale}

% --- PAQUETES ---
\usepackage[utf8]{inputenc}
\usepackage[spanish]{babel}
\usepackage{amsmath, amssymb}
\usepackage{siunitx}
\usepackage{booktabs}
\usepackage{tabularx}
\usepackage[american]{circuitikz}
\usepackage{pgfplots}
\pgfplotsset{compat=1.18}

% Definición del color darkred que requiere pgfplots
\definecolor{darkred}{rgb}{0.55,0.0,0.0}

\sisetup{output-decimal-marker = {,}, per-mode = symbol}

% --- INFORMACIÓN DEL CURSO ---
\title[Clippers y Clampers]{Circuitos Limitadores y Sujetadores}
\subtitle{Conformación No Lineal de Ondas y Acondicionamiento de Señal}
\author{M.Sc. Ing. Bernardo Quiroga Turdera}
\institute[UCB Tarija]{Circuitos Electrónicos III (IMT-221)}
\date{Semana 3: Sesión 6}

\begin{document}

\begin{frame}
  \titlepage
\end{frame}

\begin{frame}{El Reto en Mecatrónica y Biomédica (PFI)}
  \begin{block}{El Problema de Rango Dinámico}
      Sensores analógicos (piezoeléctricos, inductivos) entregan señales bipolares (ej. $\pm\qty{10}{\volt}$). \\
      Sin embargo, un ADC de microcontrolador procesa señales monopolares ($\qty{0}{\volt}$ a $\qty{+5.0}{\volt}$).
  \end{block}
  \vspace{0.2cm}
  \alert{\textbf{Si conectamos la señal directamente:}}\\
  \begin{itemize}
      \item Semiciclos negativos \textbf{destruirán la compuerta P-N} del microcontrolador.
      \item Sobretensiones positivas \textbf{quemarán térmicamente} el integrado.
  \end{itemize}
  \vspace{0.2cm}
  \textbf{Solución en 2 etapas:}\\
  1. \textit{Clampers}: Desplazar toda la onda hacia arriba (añadir offset DC).\\
  2. \textit{Clippers}: Recortar los picos destructivos sin afectar la información útil.
\end{frame}

\begin{frame}{Circuitos Recortadores (Clippers)}
  \begin{columns}
      \begin{column}{0.48\textwidth}
          \textbf{Recortador Paralelo Positivo}\\
          \begin{circuitikz}[scale=0.7, transform shape]
              \draw (0,0) to[short, o-] (1,0) to[R] (3,0) to[short, -o] (4,0);
              \draw (3,0) to[D] (3,-1.5);
              \draw (0,-1.5) to[short, o-o] (4,-1.5);
          \end{circuitikz}\\
          Si $v_i > V_\gamma \implies v_o = V_\gamma$
          
          \vspace{0.3cm}
          \textbf{Recortador Polarizado ($+V_{\text{REF}}$)}\\
          \begin{circuitikz}[scale=0.7, transform shape]
              \draw (0,0) to[short, o-] (1,0) to[R] (3,0) to[short, -o] (4,0);
              \draw (3,0) to[D] (3,-1.2) to[V, l=$V_{\text{ref}}$] (3,-2.5);
              \draw (0,-2.5) to[short, o-o] (4,-2.5);
          \end{circuitikz}\\
          Si $v_i > V_{\text{REF}} + V_\gamma \implies v_o = V_{\text{REF}} + V_\gamma$
      \end{column}
      
      \begin{column}{0.5\textwidth}
          \textbf{Curva de Transferencia ($v_o$ vs. $v_i$)}\\
          Muestra cómo la señal pasa intacta (pendiente 1) hasta chocar con el umbral.
          \vspace{0.2cm}\\
          \begin{tikzpicture}[scale=0.75]
              \begin{axis}[
                  axis lines=middle, xmin=-3, xmax=6, ymin=-3, ymax=6,
                  xlabel={$v_i$}, ylabel={$v_o$},
                  xtick={4}, xticklabels={$V_{\text{REF}} \text{+} V_\gamma$},
                  ytick={4}, yticklabels={$V_{\text{REF}} \text{+} V_\gamma$},
                  grid=both, grid style={dashed, gray!30}
              ]
              \addplot[thick, blue, domain=-3:4] {x};
              \addplot[thick, red, domain=4:6] {4};
              \end{axis}
          \end{tikzpicture}
      \end{column}
  \end{columns}
\end{frame}

\begin{frame}{Circuitos Sujetadores (Clampers): Principio Físico}
  Añade un offset de DC \textbf{sin alterar la amplitud pico a pico} ($V_{pp}$).
  \begin{columns}
      \begin{column}{0.5\textwidth}
          \textbf{Mecánica de Carga y Batería:}\\
          \begin{itemize}
              \item \textbf{Carga Transitoria:} El diodo se enciende al primer pico. $C$ se carga a $V_C = V_p - V_\gamma$.
              \item \textbf{Régimen Permanente:} El diodo se apaga. $C$ retiene la carga como una batería.
              \item \textbf{Ecuación:} $v_o(t) = v_i(t) + V_C$
          \end{itemize}
      \end{column}
      \begin{column}{0.5\textwidth}
          \textbf{Criterio de Diseño Crítico:}\\
          Para evitar que el capacitor se vacíe (\textit{droop} o decaimiento) arruinando la señal:
          \begin{block}{Condición de Descarga}
              \begin{equation*}
                  \tau = R_L \cdot C \ge 10 \cdot T
              \end{equation*}
          \end{block}
      \end{column}
  \end{columns}
\end{frame}

\begin{frame}{Problema Integrador: Acondicionamiento (Clamper + Clipper)}
  Se requiere meter una señal senoidal pura de $\pm\qty{6}{\volt}$ al ADC (\qty{0}{\volt} a \qty{5.0}{\volt}).
  \begin{center}
      \begin{circuitikz}[american, scale=0.75, transform shape]
          \draw (0,0) to[sV, l={$v_i$ ($\pm\qty{6}{\volt}$)}] (0,3) to[C, l=$C$] (2,3);
          \draw (2,0) to[D, l_=$D_1$] (2,3);
          \draw (4,3) to[R, l=$R_L$] (4,0);
          \draw (2,3) -- (4,3) -- (6,3) to[R, l=$R_1$] (8,3);
          \draw (8,3) to[D, l=$D_2$] (8,1.5) to[V, l=$V_{\text{C}}$] (8,0);
          \draw (8,3) -- (10,3) node[right]{ADC};
          \draw (0,0) -- (10,0);
          \node[above] at (1,3.5) {\textbf{ETAPA 1: CLAMPER}};
          \node[above] at (7,3.5) {\textbf{ETAPA 2: CLIPPER}};
      \end{circuitikz}
  \end{center}
  \textbf{Resultados de Diseño:}\\
  $C = \qty{2.2}{\micro\farad}$ asegura el Offset de $\qty{+5.3}{\volt}$. \\
  $V_{\text{CLIP}} = \qty{4.4}{\volt}$ asegura el recorte seguro arriba de $\qty{5.1}{\volt}$.
\end{frame}

\begin{frame}{Análisis Gráfico Dinámico: Señal en el Tiempo}
    Observen cómo en el primer semiciclo ($t < 0.5T$) la onda es idéntica a la entrada, hasta que el capacitor se carga en $\qty{-0.7}{\volt}$ e impulsa toda la onda hacia arriba para siempre.
    \begin{center}
        \begin{tikzpicture}[scale=0.9]
            \begin{axis}[
                width=12cm, height=5.5cm,
                axis lines=middle,
                xlabel={$t/T$}, ylabel={Voltaje (\unit{\volt})},
                xmin=0, xmax=2.5, ymin=-7, ymax=13,
                grid=both, grid style={dashed, gray!30},
                legend style={at={(0.98,0.05)}, anchor=south east}
            ]
            \addplot[domain=0:2.5, samples=200, gray, dashed, thick] {6*sin(360*x)};
            \addplot[domain=0:2.5, samples=300, blue, thick] {
                (x<0.5)*(6*sin(360*x)) + (x>=0.5)*(x<0.75)*(max(-0.7, 6*sin(360*x))) + (x>=0.75)*(max(-0.7, 6*sin(360*x) + 5.3))
            };
            \addplot[domain=0:2.5, samples=300, red, ultra thick] {
                min(5.1, (x<0.5)*(6*sin(360*x)) + (x>=0.5)*(x<0.75)*(max(-0.7, 6*sin(360*x))) + (x>=0.75)*(max(-0.7, 6*sin(360*x) + 5.3)))
            };
            \addplot[domain=0:2.5, darkred, dashed] {5.1};
            \addlegendentry{Entrada} \addlegendentry{Clamper ($v_{o1}$)} \addlegendentry{Protegida ($v_{ADC}$)}
            \end{axis}
        \end{tikzpicture}
    \end{center}
\end{frame}

\begin{frame}{Análisis Gráfico: Curva de Transferencia Final}
    \begin{columns}
        \begin{column}{0.4\textwidth}
            \textbf{Curva del Sistema Global}\\
            Eliminando el factor temporal, vemos qué le ocurre a cada valor instantáneo de $v_{in}$ que ingresa al ADC.
            
            \vspace{0.3cm}
            \begin{itemize}
                \item \textbf{Corte Inferior:} \qty{-0.7}{\volt}
                \item \textbf{Zona Lineal:} Pendiente 1.
                \item \textbf{Corte Superior:} \qty{5.1}{\volt}
            \end{itemize}
        \end{column}
        
        \begin{column}{0.6\textwidth}
            \begin{tikzpicture}
                \begin{axis}[
                    width=\textwidth, height=6.5cm,
                    axis lines=middle,
                    xlabel={$v_{in}$ (\unit{\volt})}, ylabel={$v_{ADC}$ (\unit{\volt})},
                    xmin=-7, xmax=7, ymin=-2, ymax=7,
                    grid=both, grid style={dashed, gray!30}
                ]
                \addplot[domain=-6:6, samples=150, purple, ultra thick] {min(5.1, max(-0.7, x + 5.3))};
                \draw[dashed, gray] (axis cs:-6, -0.7) -- (axis cs: -6, 0);
                \draw[dashed, gray] (axis cs:-0.2, 5.1) -- (axis cs: 6, 5.1);
                \end{axis}
            \end{tikzpicture}
        \end{column}
    \end{columns}
\end{frame}

\end{document}